\section{Explicit finite data for the complete \texorpdfstring{$c=3$}{c=3} slice}\label{sec:certificates}
This section prints every finite witness used by the periodic construction and the finite bridge.
Thus those parts of the existence proof are logically independent of the accompanying programs:
the programs merely regenerate and check the data below.  All strings are read from left to right.
Throughout this section \(n=3m\), \(\kappa=(n+1)/2\), and all phases are taken modulo \(24\).
The symbols \(L\), \(M\), and \(U\) in the boundary table refer to the lower endpoint,
the \(\kappa\)-window, and the upper endpoint, respectively.  The bulk blocks are denoted
by \(B_1,B_2,B_3\).  Thus a certificate is read in three steps: decode the offset word,
lift it with the formula in Section~\ref{sec:construction}, and then apply the displayed orientation
and sign strings.

\subsection{Offset words and boundary tables}\label{sec:offset-data}
We encode each zero-sum offset vector by the following alphabet.  Every listed symbol belongs
to \(\{-2,-1,0,1,2\}^3\) and has coordinate sum zero; consequently the row-sum condition
is built into the encoding rather than checked by a separate search.
\begingroup\scriptsize
\[
\begin{array}{c|ccccccc}
\text{symbol}&0&1&2&3&4&5&6\\\hline
\text{offset}&(-2,0,2)&(-1,-1,2)&(-1,0,1)&(0,-2,2)&(0,-1,1)&(0,0,0)&(0,1,-1)\\
\end{array}
\]
\endgroup
\begingroup\scriptsize
\[
\begin{array}{c|cccccc}
\text{symbol}&7&8&9&A&B&C\\\hline
\text{offset}&(1,-2,1)&(1,-1,0)&(1,0,-1)&(2,-2,0)&(2,-1,-1)&(2,0,-2)\\
\end{array}
\]
\endgroup
For a word indexed by \(0,1,\ldots,23\), the first printed symbol is the value at index \(0\).
The complete bulk data reduce to the following two master words:
\[
\begin{array}{c|c}
W_1&\mathtt{8AA2375331378AA537533437}\\
W_2&\mathtt{4331347AA5344334347AA234}
\end{array}
\]
For $n\equiv9\pmod{24}$ use $(W_1,W_2)$ with phases $(6,12)$.  For
$n\equiv21\pmod{24}$ use phases $(10,13)$ and
\begin{equation}\label{eq:c3-word-shift}
 \widehat W_1[p]=W_1[p+16\bmod24],\qquad
 \widehat W_2[p]=W_2[p+13\bmod24].
\end{equation}
Thus the formerly separate words for the second class are exact rotations of the displayed words;
after the phases are inserted, passage between the two classes is a half-period shift by $12$ in
both bulk layers.
For each boundary window, the symbols are listed in increasing order of the indicated relative
index; no reversal of a printed word is intended.
Here \(L,M,U\) mean the windows at \(1,\kappa,n\), respectively.
\[
\begin{array}{c|c|c|l}
n\bmod48&\text{window}&\text{indices}&\text{offset word}\\\hline
9&L&1:20&\mathtt{AAA57AACA8A5331378AA}\\
9&M&-20:20&\mathtt{2375331378AA537A534378A81344331347AA53443}\\
9&U&-20:0&\mathtt{3347137A53432347AAA80}\\
33&L&1:20&\mathtt{57AB88A647A5331378AA}\\
33&M&-20:20&\mathtt{5375334378AA237A530343433344334347AA23443}\\
33&U&-20:0&\mathtt{33431A7A53478547AAAB8}\\
21&L&1:20&\mathtt{A988AAA8AA75334378AA}\\
21&M&-20:20&\mathtt{4378AA2375331378AAA57AA8A547AA2344331347A}\\
21&U&-20:0&\mathtt{3343437A5343334717570}\\
45&L&1:20&\mathtt{A81347A5337533437833}\\
45&M&-20:20&\mathtt{1378AA5375334378AA2375331347A53344334347A}\\
45&U&-20:0&\mathtt{33431337534344333447A}\\
\end{array}
\]
For an independent support audit, fix one of the $16$ admissible periodic states described
below and a residue \(u\in\{1,\ldots,72\}\).  Write $n_0=e+24\tau$, so that
$n=n_0+288Q$ with $Q\ge0$.  Decoding the two bulk words and the corresponding boundary
table turns every family of values congruent to $u\pmod{72}$ into an integer interval
whose endpoints are affine functions of \(Q\).  For each of the $16\times72$ state--residue
pairs, the displayed boundary word gives the endpoint relation at \(Q=0\), while the period
words give the same endpoint slopes for every \(Q\ge0\).  Hence the intervals remain consecutive,
begin at \(0\), and end at
\(12Q+\lfloor(3n_0-u)/72\rfloor\), exactly as required in
Lemma~\ref{lem:rule}.  This is a finite comparison of printed affine expressions, not a
computer-generated existence assertion.

\subsection{The 24 stored orientation patterns for 48 admissible affine states}\label{sec:orientation-data}
For a row index \(j\), let its category be
\[
\gamma(j)=\begin{cases}
(L,j),&j\le20,\\
(M,j-m),&|j-m|\le12,\\
(K,j-\kappa),&|j-\kappa|\le20,\\
(T,j-n),&j\ge n-20,\\
(B_1,j\bmod24),&j<m-12,\\
(B_2,j\bmod24),&m+12<j<\kappa-20,\\
(B_3,j\bmod24),&\text{otherwise}.
\end{cases}
\]
The cases are applied in the displayed order.  For any admissible parameter using the periodic rule,
put \(e=n\bmod24\) (with the representatives \(e=9,21\) in the periodic range),
\(t=(n-e)/24\), and select the stored row with
\(s=15+((t-15)\bmod6)\).  Coprimality gives
\[
s\in\{16,17,19,20\}\quad(e=9),\qquad
s\in\{15,17,18,20\}\quad(e=21).
\]
This also specifies the ten finite-bridge cases below.  For those bridge values, \(s\) is only
the table-selection residue; the affine \(\tau\)-range starts at \(15\).  For the infinite proof, an admissible affine
group state is \((e,\tau,r)\), where
\(n=e+24(\tau+12Q)\), \(e\in\{9,21\}\), \(15\le\tau\le26\), \(Q\ge0\), and
the retained indices satisfy \(j\equiv r\pmod3\).  The stored row \((e,s,r)\),
\(15\le s\le20\), is used for both \(\tau=s\) and \(\tau=s+6\).  For either value of
\(\tau\), let \(\mathcal A_{e,\tau,r}\) be the union of its categories for \(Q=0,1,\ldots,6\);
the literal category set is the same for the two values represented by a stored row.
The restrictions above leave $8$ stored parameter states and hence $24=8\cdot3$ stored rows; they
cover the $48=16\cdot3$ admissible affine group states.  Here \(r\in\{0,1,2\}\) is the group
index, \(p=j\bmod24\in\{0,\ldots,23\}\), and \(z\in\{1,2,3\}\) is the bulk zone.  The
category-to-label dictionary used by the printed strings is
\[
L\longleftrightarrow\mathtt{('bnd','low',q)},\quad
M\longleftrightarrow\mathtt{('bnd','m',q)},\quad
K\longleftrightarrow\mathtt{('bnd','k',q)},\quad
T\longleftrightarrow\mathtt{('bnd','top',q)}.
\]
The letters \(M\) and \(K\) in the category list are deliberately distinct: here \(M\) denotes
the window centred at the integer \(m\), whereas \(K\) denotes the window centred at \(\kappa\).
This category notation is independent of the \(L,M,U\) labels used for the offset-table display
above; the literal labels in the sign strings are the authoritative ones for decoding.
Thus a switch \(K:q\) refers to the literal key \(\mathtt{('bnd','k',q)}\), while a bulk key
has the form \(\mathtt{('bulk',z,p)}\) with \(z\in\{1,2,3\}\) and \(p\in\{0,\ldots,23\}\).
For encoding,
write the labels literally as \(\mathtt{('bnd','k',q)}\), \(\mathtt{('bnd','low',q)}\),
\(\mathtt{('bnd','m',q)}\), \(\mathtt{('bnd','top',q)}\), and
\(\mathtt{('bulk',z,p)}\), and order these literal strings lexicographically.  Each hexadecimal digit below
expands to four binary digits; retain the first \(\ell\) bits, with \(1=+1\) and \(0=-1\).
The columns \(C,D\) give respectively the signs of \(\mu_j\) and \(|w_j|\); these are sign-string
columns and are unrelated to the column-sum vectors \(C_i\) in the finite-bridge table below.
Both columns use the same ordered list of literal keys, whose length is \(\ell\).  A switch
\(X:q;a\) or \(X:q;b\) changes the unique row in category \((X,q)\) from mode \(c\)
to the displayed mode; ``--'' means that every row uses mode \(c\).
\begin{longtable}{c@{\,}c@{\,}c|r|c|l|l}
\hline
e&s&r&$\ell$&\text{switch}&C&D\\\hline
\endfirsthead
\hline e&s&r&$\ell$&\text{switch}&C&D\\\hline
\endhead
9&16&0&59&--&\texttt{FE7FAF801E1E1E0}&\texttt{BFFFFC481FF51C0}\\
9&16&1&60&\(K:-1;b\)&\texttt{FF3B87F00F0B8F0}&\texttt{FFFEE80404DB8F0}\\
9&16&2&60&--&\texttt{FEDFAF00069F0F0}&\texttt{FF7FF0140AA78F0}\\
9&17&0&59&\(K:-11;a\)&\texttt{7FFE55001E1E1E0}&\texttt{BEFBF0081FF51C0}\\
9&17&1&61&--&\texttt{FDFE07C0054F8780}&\texttt{9FFFFC0202CDC780}\\
9&17&2&59&--&\texttt{F7FD0F800D3E1E0}&\texttt{FAFFF808154F1E0}\\
9&19&0&59&\(K:-11;a\)&\texttt{7FFD3C001E1E1E0}&\texttt{F6F7F0081FF51C0}\\
9&19&1&60&--&\texttt{FFABC7C00A9F0F0}&\texttt{FFFF7A04059B8F0}\\
9&19&2&60&--&\texttt{DDF7AF00069F0F0}&\texttt{FBDFF4040AA78F0}\\
9&20&0&59&--&\texttt{FFEBCF801E1E1E0}&\texttt{DFFFFC081FF51C0}\\
9&20&1&61&\(K:-1;b\)&\texttt{FF3B87F00785C780}&\texttt{FFFFFC00026DC780}\\
9&20&2&59&--&\texttt{FFF8AF200D3E1E0}&\texttt{FF7FF028154F1E0}\\
21&15&0&59&--&\texttt{FFDFF000153E1E0}&\texttt{FFBFFE481FBBFC0}\\
21&15&1&61&--&\texttt{FFF847E0054F8780}&\texttt{DFF7FBA202CDC780}\\
21&15&2&59&--&\texttt{BFEF6F800D3E1E0}&\texttt{FFEFFF58154F1E0}\\
21&17&0&59&--&\texttt{EF7BE800153E1E0}&\texttt{FBFFF8081FBBFC0}\\
21&17&1&60&--&\texttt{FEFC87C00A9F0F0}&\texttt{BF7FFE84059B8F0}\\
21&17&2&60&--&\texttt{7FF76F80069F0F0}&\texttt{FFF7FF440AA78F0}\\
21&18&0&59&--&\texttt{FFF52E00153E1E0}&\texttt{BFDFF4081FBBFC0}\\
21&18&1&61&\(K:-10;a\)&\texttt{BFAD87E0054F8780}&\texttt{6F77F80202CDC780}\\
21&18&2&59&\(K:-12;a\)&\texttt{FDDF8F800D3E1E0}&\texttt{EFCFF008154F1E0}\\
21&20&0&59&--&\texttt{FFEC5800153E1E0}&\texttt{7DFFF0081FBBFC0}\\
21&20&1&60&\(K:-10;a\)&\texttt{FEAF17C00A9F0F0}&\texttt{FFAEF804059B8F0}\\
21&20&2&60&\(K:-12;a\)&\texttt{7DFE4F80069F0F0}&\texttt{EFAFF8040AA78F0}\\
\hline
\end{longtable}
For direct checking, substitute the offsets from Section~\ref{sec:offset-data} in the
closed lift formula.  For each admissible stored row, for the two values \(\tau=s,s+6\), and in each category,
the sums of \(\mu_j\) and \(|w_j|\) are polynomials in \(Q\) of degree at most two.  The two decoded
sign rows make their signed totals zero at \(Q=0,1,\ldots,6\) for both values of \(\tau\);
seven zeros therefore force both polynomials to vanish identically in all \(48\) admissible states.
This is a finite arithmetic verification, and no search or program is part of the logical implication.
\subsection{The ten-case finite bridge}\label{sec:bridge-data}
The periodic support rule is used from \(m\ge123\) onward.  The following ten admissible values
of \(m\) bridge the remaining interval \(67\le m\le119\).  Let \(S_m\) be the set of absolute
values produced and let
\(C_i\) be the three-component column-sum vector of the \(i\)-th array.  Direct substitution
in the printed rules gives the complete audit table below.  Here
\(C_i=(C_{i,1},C_{i,2},C_{i,3})\).  Since every row sum is zero by construction,
\(S_m\subseteq[1,9m]\); the equalities \(\min S_m=1\), \(\max S_m=9m\), and
\(|S_m|=9m\) show that \(S_m\) consists of \(9m\) distinct integers in an interval of
length \(9m\), and hence \(S_m=[1,9m]\).  The last column gives \(C_i=\mathbf0\) for
each \(i\), which is exactly the vanishing of all nine column sums.
\[
\begin{array}{c|c|c|c|c|c}
m&n&n\bmod48&s&(\min S_m,\max S_m,|S_m|)&(C_0,C_1,C_2)\\\hline
67&201&9&20&(1,603,603)&(\mathbf{0},\mathbf{0},\mathbf{0})\\
71&213&21&20&(1,639,639)&(\mathbf{0},\mathbf{0},\mathbf{0})\\
79&237&45&15&(1,711,711)&(\mathbf{0},\mathbf{0},\mathbf{0})\\
83&249&9&16&(1,747,747)&(\mathbf{0},\mathbf{0},\mathbf{0})\\
91&273&33&17&(1,819,819)&(\mathbf{0},\mathbf{0},\mathbf{0})\\
95&285&45&17&(1,855,855)&(\mathbf{0},\mathbf{0},\mathbf{0})\\
103&309&21&18&(1,927,927)&(\mathbf{0},\mathbf{0},\mathbf{0})\\
107&321&33&19&(1,963,963)&(\mathbf{0},\mathbf{0},\mathbf{0})\\
115&345&9&20&(1,1035,1035)&(\mathbf{0},\mathbf{0},\mathbf{0})\\
119&357&21&20&(1,1071,1071)&(\mathbf{0},\mathbf{0},\mathbf{0})\\
\end{array}
\]
